\chapter{Introducción}

La industria del videojuego se ha consolidado durante las últimas décadas como uno de los sectores tecnológicos de mayor crecimiento, superando en volumen de negocio a otras industrias culturales tradicionales \cite{GAMES_REPORT}. Desde la perspectiva de la ingeniería del software, el desarrollo de videojuegos plantea desafíos técnicos relevantes relacionados con la arquitectura de sistemas en tiempo real, la implementación de algoritmos de inteligencia artificial, la generación procedural de contenido y el diseño de interfaces interactivas. Este contexto convierte el desarrollo de videojuegos en un ámbito adecuado para la realización de un Trabajo de Fin de Grado en Ingeniería Informática, al permitir la aplicación práctica de conocimientos adquiridos durante la formación académica.

\section{Contexto y relevancia del problema}

Los videojuegos pertenecientes al subgénero \textit{rogue-like} presentan una serie de características técnicas que los convierten en un objeto de estudio de especial interés. Entre ellas destacan la generación procedural de niveles, la ausencia de persistencia del progreso entre partidas, la progresión limitada a cada ejecución del juego y una elevada rejugabilidad derivada de la variabilidad del contenido generado. Estas propiedades obligan a diseñar sistemas capaces de producir escenarios jugables coherentes y equilibrados de forma automática, sin intervención manual directa.

En este contexto se plantea el siguiente problema técnico: cómo diseñar e implementar una arquitectura modular que permita la generación automática de niveles jugables, la gestión concurrente de múltiples subsistemas —como la inteligencia artificial, la física, la interfaz de usuario o el control de estados— y la parametrización del contenido de manera escalable y mantenible.

La actualidad de este problema se pone de manifiesto en la proliferación de títulos comerciales pertenecientes al género en los últimos años, tales como Hades, Dead Cells, The Binding of Isaac o Enter the Gungeon. Desde un punto de vista académico, el análisis de estos sistemas permite abordar conceptos relacionados con la arquitectura de software, el uso de patrones de diseño, los algoritmos de generación procedural y la separación entre lógica de ejecución y datos configurables.

\section{Propósito del trabajo}

El propósito del presente Trabajo de Fin de Grado consiste en el diseño, la implementación y la documentación de un videojuego funcional perteneciente al subgénero \textit{rogue-like}, utilizando Unity como motor gráfico \cite{REF_UNITY} y C\# como lenguaje de programación. El objetivo principal es demostrar la aplicabilidad de principios propios de la ingeniería del software en un contexto práctico, prestando especial atención a la modularidad del sistema, su escalabilidad y la automatización del flujo de trabajo mediante herramientas de desarrollo.

El proyecto aborda, de forma integrada, tres áreas de conocimiento fundamentales. En primer lugar, la arquitectura de software, mediante la aplicación de patrones de diseño como \textit{Singleton}, \textit{Observer} y \textit{Strategy}, así como el uso de \textit{ScriptableObjects} para separar la lógica de juego de los datos configurables en aquellos sistema donde resulta adecuado. En segundo lugar, los algoritmos y estructuras de datos, a través de la implementación de sistemas de generación procedural de salas, la gestión ponderada de enemigos y el control del flujo del juego mediante máquinas de estados. Por último, la ingeniería de herramientas, mediante el desarrollo de utilidades dentro del editor de Unity orientadas a automatizar tareas repetitivas y a reducir la probabilidad de error humano durante la creación de contenido.

\section{Descripción del proyecto}

El videojuego desarrollado, denominado Ascension, es un \textit{rogue-like} bidimensional con vista cenital, en el que el jugador avanza a través de una secuencia de salas generadas de forma procedural, enfrentándose a enemigos hasta alcanzar la condición de victoria establecida. Dicha condición consiste en la derrota de tres jefes que aparecen de forma periódica cada cinco salas.

El sistema implementado incluye un flujo de juego completo que abarca el menú principal, la selección de clase entre tres personajes con estadísticas diferenciadas, la gestión de estados globales mediante un \textit{GameManager} y las pantallas de victoria o derrota. Asimismo, se ha desarrollado un sistema de generación procedural de salas basado en \textit{tilemaps}, con muros perimetrales, elementos de transición entre salas basados en \textit{tiles} interactivos (escaleras).

El sistema de combate contempla dos tipos de armas, cuerpo a cuerpo y a distancia, con proyectiles físicos, detección de colisiones y un sistema configurable de armas iniciales que está preparado para extensión futura. La inteligencia artificial de los enemigos incluye tres comportamientos diferenciados —persecución, ataque a distancia y aproximación por saltos—, además de enemigos especiales con atributos incrementados que actúan como jefes de sala. La gestión de enemigos se realiza mediante generación ponderada por coste, limpieza automática de salas y control de interacción con escaleras en función del estado del combate.

El proyecto incorpora, además, una interfaz de usuario completa con visualización de la salud del jugador, mensajes contextuales y una tabla de puntuaciones persistente. De forma complementaria, se han desarrollado herramientas de automatización dentro del editor de Unity destinadas a la configuración automática de salas, la generación de salas, la gestión gráfica de bases de datos de enemigos y armas, y la aplicación de efectos sobre \textit{tiles}. El sistema no incluye gestión de audio ni un controlador genérico de objetos recolectables, al considerarse estos aspectos de menor prioridad para los objetivos académicos del trabajo.

\section{Objetivos}

El objetivo general del proyecto es desarrollar un videojuego \textit{rogue-like} funcional que permita evidenciar la aplicación de principios de ingeniería del software en el ámbito del desarrollo de videojuegos, con especial énfasis en la arquitectura modular, la generación procedural de contenido y la automatización de flujos de trabajo.

De este objetivo general se derivan los siguientes objetivos específicos: diseñar e implementar una arquitectura modular basada en patrones de diseño que facilite la separación de responsabilidades entre subsistemas; desarrollar un algoritmo de generación procedural de salas que garantice la jugabilidad y la coherencia espacial del entorno; implementar distintos comportamientos de inteligencia artificial para enemigos y enemigos especiales con atributos incrementados que actúan como jefes de sala; crear un sistema de armas con generación procedural de estadísticas y mecánicas de combate cuerpo a cuerpo y a distancia; desarrollar herramientas de automatización dentro del editor de Unity para agilizar la creación y validación de contenido; implementar un sistema de persistencia de puntuaciones mediante \textit{PlayerPrefs} con visualización de un ranking en el menú principal; y documentar el proceso de desarrollo, las decisiones técnicas adoptadas y las limitaciones encontradas.

El presente documento expone el análisis, el diseño y la implementación detallada de estos sistemas, así como las conclusiones derivadas del proceso de desarrollo.
